\section{Related Work}

\paragraph{Positional encoding design space.}
Transformers rely on positional structure to recover order information \citep{vaswani2017attention}. The design space includes relative attention biases \citep{shaw2018self, raffel2020exploring}, additive linear biases \citep{press2022alibi}, kernelized positional embeddings \citep{chi2022kerple}, and rotary position embedding (RoPE) \citep{su2024roformer}. Within RoPE, prior work operates on two axes: the base frequency $b$ (NTK-aware scaling, YaRN) and inference-time position rescaling. A third axis---the allocation of frequencies \emph{across channels}---has remained at the geometric default since RoFormer. Resonance RoPE~\citep{wang2024resonance} is the closest precedent: it identifies critical frequencies that should not be perturbed during interpolation, implicitly acknowledging that not all channels are equal. EVQ addresses the complementary question: given that channels differ in importance, what is the optimal \emph{training-time} allocation?

\paragraph{Inference-time and post-training context extension.}
These methods fall into three categories. \emph{Pure inference-time rescaling}: positional interpolation~\citep{chen2024position} and YaRN~\citep{peng2024yarn} rescale RoPE positions on a frozen checkpoint; SelfExtend~\citep{jin2024selfextend} uses grouped attention without fine-tuning. \emph{Search-based or progressive adaptation}: LongRoPE~\citep{ding2024longrope} and PoSE~\citep{zhu2024pose} use progressive or skip-wise strategies; LongRoPE2~\citep{shang2025longrope2} searches for per-channel scaling factors via evolution on a frozen checkpoint. \emph{Trained extrapolation modules}: CLEX~\citep{chen2024clex} trains a continuous length-extrapolation module. All three categories preserve the pretrained checkpoint and modify behavior at or after inference time. EVQ is orthogonal: it changes the training-time frequency layout itself. While LongRoPE2 and EVQ both operate on per-channel frequency structure, they target different design axes: EVQ modifies the layout \emph{before training}, whereas LongRoPE2 searches for rescaling \emph{after training}. This complementarity is demonstrated empirically: EVQ+YaRN substantially outperforms Geo+YaRN (Table~\ref{tab:evq-yarn}), confirming that a better training-time substrate amplifies inference-time methods rather than competing with them.

\paragraph{Learnable and adaptive positional encoding.}
Another line of work introduces learnable or adaptive positional mechanisms: FIRE~\citep{li2024fire} learns a continuous interpolation function, and DAPE~\citep{zheng2024dape} adapts positional encoding to data statistics. These methods are the closest comparison point for our PE-dominant experiments because they ask whether adaptive positional structure can improve extrapolation. Our distinction is that EVQ is derived in closed form rather than learned through additional positional parameters, which lets us separate the question of \emph{what frequency allocation is optimal} from the question of \emph{whether an optimizer can discover it from in-distribution loss alone}.

\paragraph{Video positional encoding.}
Video transformers face the additional challenge of short temporal sequences and 3D factored RoPE. VideoRoPE~\citep{videorope2025} is the closest video positional-design work, identifying sequence-length-dependent frequency allocation as a key design axis for video generation. RIFLEx~\citep{zhao2025riflex} is an inference-time frequency correction method for video diffusion transformers. FreeNoise~\citep{qiu2024freenoise} addresses temporal coherence via noise rescheduling and is adjacent but not a positional-encoding method. EVQ is complementary to all three: it optimizes the training-time allocation itself, and preliminary video DiT experiments (2~seeds) confirm that the collision-reduction mechanism extends beyond text.

\paragraph{Evaluation for long-context behavior.}
Long-context evaluation spans synthetic retrieval diagnostics and broader task suites \citep{bai2024longbench, hsieh2024ruler}. Prior work shows that long contexts can fail even when relevant information is nominally present \citep{liu2024lost}, and that positional encoding modifications can mitigate such failures \citep{zhang2024found}. We use both PE-dominant settings (isolating frequency-allocation effects) and retrieval-heavy evaluation (testing systems-level claims), treating incomplete larger-scale evidence conservatively.
